%% =====================================================================
%%  common/merda.tex — MER-DA の資料に共通の部品（上端の帯，補足のページを指す命令，色，図の部品）
%%  2026-10-27 に ~/slides/decks/MER-DA/2026-10-22-6つの問いの整理/main.tex の前置きから写した．
%%  以後はこのファイルを正本とし，~/slides とは同期しない．
%%  読み込み：各資料の main.tex で %% 2026-10-27 に ~/slides/common/ から写した（MER-DA を自己完結させるため）．以後は MER-DA/slides/common/ のこのファイルを正本とし，~/slides とは同期しない．
%% =====================================================================
%%  common/preamble.tex — 全スライド共通のプリアンブル
%%  エンジン : LuaLaTeX（日本語は luatexja）
%%  読み込み : 各 deck の main.tex から %% 2026-10-27 に ~/slides/common/ から写した（MER-DA を自己完結させるため）．以後は MER-DA/slides/common/ のこのファイルを正本とし，~/slides とは同期しない．
%% =====================================================================
%%  common/preamble.tex — 全スライド共通のプリアンブル
%%  エンジン : LuaLaTeX（日本語は luatexja）
%%  読み込み : 各 deck の main.tex から %% 2026-10-27 に ~/slides/common/ から写した（MER-DA を自己完結させるため）．以後は MER-DA/slides/common/ のこのファイルを正本とし，~/slides とは同期しない．
%% =====================================================================
%%  common/preamble.tex — 全スライド共通のプリアンブル
%%  エンジン : LuaLaTeX（日本語は luatexja）
%%  読み込み : 各 deck の main.tex から \input{preamble}
%%             （common/ は .latexmkrc が TEXINPUTS に通す）
%%
%%  読み込み順には意味がある．並べ替えないこと．
%%    1. luatexja        … 和文組版の土台．最初に読む
%%    2. theme.tex       … \usetheme{metropolis}．内部で fontspec を読む
%%    3. luatexja-preset … fontspec より後に和文フォントを確定させる
%%    4. 一般パッケージ / マクロ
%% =====================================================================

% ---------- 1. 日本語（土台）------------------------------------------
\usepackage{luatexja}

% ---------- 2. テーマ（metropolis）------------------------------------
\input{theme}

% ---------- 3. 和文フォント -------------------------------------------
% Harano Aji は texlive-lang-japanese 同梱なので環境依存が少ない．
% Noto CJK を使う場合は [haranoaji] を [noto-otc] に差し替える．
\usepackage[haranoaji,deluxe]{luatexja-preset}
% スライドの既定は明朝ではなくゴシック（プロジェクタでの可読性）
\renewcommand{\kanjifamilydefault}{\gtdefault}
% 和欧文間の空きをやや詰める（スライドは行長が短いため）
\ltjsetparameter{xkanjiskip=0.20\zw plus 0.1\zw minus 0.05\zw}

% ---------- 4. 数式 ----------------------------------------------------
\usepackage{amsmath,amssymb}
\usepackage{mathtools}
\usepackage{bm}
\usepackage{siunitx}
\sisetup{detect-all, per-mode=symbol, group-minimum-digits=4}

% ---------- 5. 図表 ----------------------------------------------------
\usepackage{graphicx}
\usepackage{booktabs}
\usepackage{subcaption}
\usepackage{adjustbox}
\graphicspath{{figures/}{./}}

% ---------- 6. 作図 ----------------------------------------------------
\usepackage{tikz}
\usetikzlibrary{arrows.meta,positioning,calc,fit,backgrounds,shapes.geometric}
\usepackage{pgfplots}
\pgfplotsset{compat=1.18}
% pgfplots の既定を「論文の figure」ではなく「スライド」向けに寄せる
\pgfplotsset{
  every axis/.append style={
    tick label style={font=\scriptsize},
    label style={font=\small},
    legend style={font=\scriptsize, draw=none, fill=none},
    line width=0.9pt, tick style={line width=0.6pt},
  },
}

% ---------- 7. 参考文献 ------------------------------------------------
% biblatex/biber は未導入のため BibTeX + natbib を使う
\usepackage[numbers,sort&compress]{natbib}
\bibliographystyle{abbrvnat}
\setbeamertemplate{bibliography item}[text]

% ---------- 8. 体裁 ----------------------------------------------------
\usepackage{appendixnumberbeamer}   % \appendix 以降をページ番号の分母から除く
\usepackage{xcolor}
\usepackage{microtype}

% 図表キャプションは番号なし（スライド中で番号参照しないため）
\setbeamertemplate{caption}{\insertcaption}
\setbeamerfont{caption}{size=\scriptsize}

% リンク色（metropolis の配色に合わせる）
% 目次とセクション扉は内部リンクなので linkcolor が効く．
% 既定のオレンジは白地で読みにくいため，地の文と同じ色にする．
\hypersetup{colorlinks=true, allcolors=., urlcolor=mAccentBlue,
            linkcolor=mDarkTeal, citecolor=mAccentBlue}

% ---------- 9. マクロ --------------------------------------------------
\input{macros}

%%             （common/ は .latexmkrc が TEXINPUTS に通す）
%%
%%  読み込み順には意味がある．並べ替えないこと．
%%    1. luatexja        … 和文組版の土台．最初に読む
%%    2. theme.tex       … \usetheme{metropolis}．内部で fontspec を読む
%%    3. luatexja-preset … fontspec より後に和文フォントを確定させる
%%    4. 一般パッケージ / マクロ
%% =====================================================================

% ---------- 1. 日本語（土台）------------------------------------------
\usepackage{luatexja}

% ---------- 2. テーマ（metropolis）------------------------------------
%% 2026-10-27 に ~/slides/common/ から写した（MER-DA を自己完結させるため）．以後は MER-DA/slides/common/ のこのファイルを正本とし，~/slides とは同期しない．
%% =====================================================================
%%  common/theme.tex — 見た目の設定（metropolis テーマ）
%%
%%  metropolis を選ぶ理由:
%%    - 装飾が最小限で，1枚あたりの情報量（図・数式）を最大化できる
%%    - フレームタイトルが大きく，「主張文をタイトルにする」書き方と相性が良い
%%    - プロジェクタで潰れにくい高コントラストの既定配色
%%    - 機械学習・物理系の国際会議発表で広く使われており，見慣れた形式
%% =====================================================================

\usetheme[
  progressbar=frametitle,   % フレームタイトル下に進捗バー
  numbering=fraction,       % 「7/23」形式のページ番号
  block=fill,               % block 環境に淡い背景を敷く
  sectionpage=progressbar,  % セクション扉に進捗バー
  subsectionpage=none,
  titleformat=regular,      % スモールキャップスにしない（和文と混ぜるため）
]{metropolis}

% ---------- 欧文フォント ----------------------------------------------
% metropolis 本来の指定は Fira Sans だが，本環境には未導入（警告が出るだけで
% ビルドは通る）．導入済みの Lato を代替に使う．Fira を入れた場合はこの
% ブロックを丸ごとコメントアウトすればテーマ既定に戻る．
\IfFontExistsTF{Lato}{\setsansfont{Lato}}{}
\IfFontExistsTF{DejaVu Sans Mono}{\setmonofont[Scale=0.85]{DejaVu Sans Mono}}{}

% ---------- 配色 ------------------------------------------------------
% metropolis の既定 (mDarkTeal / mLightBrown) を土台に，
% 「強調」「良い結果」「悪い結果」を区別できる色を足す．
\definecolor{mDarkTeal}{HTML}{23373B}    % 地の文・見出し
\definecolor{mLightBrown}{HTML}{EB811B}  % アクセント（alert）
\definecolor{mMediumBrown}{HTML}{C87A2F}
\definecolor{mAccentBlue}{HTML}{2A6FA8}  % 提案手法・注目対象
\definecolor{mGood}{HTML}{1B7F5C}        % 改善・良い結果
\definecolor{mBad}{HTML}{B3312B}         % 劣化・悪い結果
\definecolor{mMuted}{HTML}{7A8A8E}       % 補足・出典

% metropolis 既定の mLightBrown（オレンジ）は，白地に細い文字だと視認性が低い．
% alerted text はセクション扉と目次の文字色にも使われるので，青に置き換える．
\setbeamercolor{alerted text}{fg=mAccentBlue}
\setbeamercolor{progress bar}{fg=mAccentBlue}
\setbeamercolor{title separator}{fg=mAccentBlue}
\setbeamercolor{progress bar in head/foot}{fg=mAccentBlue}
\setbeamercolor{progress bar in section page}{fg=mAccentBlue}
% ---------- 余白・行間 ------------------------------------------------
\setbeamersize{text margin left=0.9em, text margin right=0.9em}
\setbeamertemplate{itemize items}[circle]
\setlength{\parskip}{0.4em}

% 箇条書きの階層ごとの間隔を詰める（1枚に収める調整用）
\setlength{\leftmargini}{1.2em}
\setlength{\leftmarginii}{1.0em}

% ---------- フッタに所属を出す ----------------------------------------
% \institute{} を設定していれば下部に小さく表示される
\setbeamertemplate{frame footer}{\insertshortinstitute}
\setbeamerfont{page number in head/foot}{size=\tiny}
\setbeamercolor{frame footer}{fg=mMuted}


% ---------- 3. 和文フォント -------------------------------------------
% Harano Aji は texlive-lang-japanese 同梱なので環境依存が少ない．
% Noto CJK を使う場合は [haranoaji] を [noto-otc] に差し替える．
\usepackage[haranoaji,deluxe]{luatexja-preset}
% スライドの既定は明朝ではなくゴシック（プロジェクタでの可読性）
\renewcommand{\kanjifamilydefault}{\gtdefault}
% 和欧文間の空きをやや詰める（スライドは行長が短いため）
\ltjsetparameter{xkanjiskip=0.20\zw plus 0.1\zw minus 0.05\zw}

% ---------- 4. 数式 ----------------------------------------------------
\usepackage{amsmath,amssymb}
\usepackage{mathtools}
\usepackage{bm}
\usepackage{siunitx}
\sisetup{detect-all, per-mode=symbol, group-minimum-digits=4}

% ---------- 5. 図表 ----------------------------------------------------
\usepackage{graphicx}
\usepackage{booktabs}
\usepackage{subcaption}
\usepackage{adjustbox}
\graphicspath{{figures/}{./}}

% ---------- 6. 作図 ----------------------------------------------------
\usepackage{tikz}
\usetikzlibrary{arrows.meta,positioning,calc,fit,backgrounds,shapes.geometric}
\usepackage{pgfplots}
\pgfplotsset{compat=1.18}
% pgfplots の既定を「論文の figure」ではなく「スライド」向けに寄せる
\pgfplotsset{
  every axis/.append style={
    tick label style={font=\scriptsize},
    label style={font=\small},
    legend style={font=\scriptsize, draw=none, fill=none},
    line width=0.9pt, tick style={line width=0.6pt},
  },
}

% ---------- 7. 参考文献 ------------------------------------------------
% biblatex/biber は未導入のため BibTeX + natbib を使う
\usepackage[numbers,sort&compress]{natbib}
\bibliographystyle{abbrvnat}
\setbeamertemplate{bibliography item}[text]

% ---------- 8. 体裁 ----------------------------------------------------
\usepackage{appendixnumberbeamer}   % \appendix 以降をページ番号の分母から除く
\usepackage{xcolor}
\usepackage{microtype}

% 図表キャプションは番号なし（スライド中で番号参照しないため）
\setbeamertemplate{caption}{\insertcaption}
\setbeamerfont{caption}{size=\scriptsize}

% リンク色（metropolis の配色に合わせる）
% 目次とセクション扉は内部リンクなので linkcolor が効く．
% 既定のオレンジは白地で読みにくいため，地の文と同じ色にする．
\hypersetup{colorlinks=true, allcolors=., urlcolor=mAccentBlue,
            linkcolor=mDarkTeal, citecolor=mAccentBlue}

% ---------- 9. マクロ --------------------------------------------------
%% 2026-10-27 に ~/slides/common/ から写した（MER-DA を自己完結させるため）．以後は MER-DA/slides/common/ のこのファイルを正本とし，~/slides とは同期しない．
%% =====================================================================
%%  common/macros.tex — 全スライド共通のマクロ
%%  ここに置くのは「どの発表でも使うもの」だけ．
%%  その発表固有のマクロは deck 側の main.tex に書く．
%% =====================================================================

% ---------- 強調 ------------------------------------------------------
\newcommand{\key}[1]{\textcolor{mAccentBlue}{\textbf{#1}}}  % 提案手法・注目対象
\newcommand{\good}[1]{\textcolor{mGood}{#1}}                % 改善・良い結果
\newcommand{\bad}[1]{\textcolor{mBad}{#1}}                  % 劣化・悪い結果
\newcommand{\muted}[1]{\textcolor{mMuted}{#1}}              % 補足

% 出典表示（右下に小さく．スライドの引用はこれで足りることが多い）
\newcommand{\source}[1]{%
  \vspace{-0.3em}\hfill{\tiny\color{mMuted}[#1]}\par}

% 「ひとことメッセージ」帯：そのスライドで言いたいことを1行で置く
\newcommand{\takeaway}[1]{%
  \vspace{0.4em}%
  \begin{beamercolorbox}[wd=\textwidth,sep=0.5em,rounded=true]{block body}%
    \textbf{#1}%
  \end{beamercolorbox}}

% ---------- 数式（機械学習でよく使うもの）-----------------------------
\newcommand{\R}{\mathbb{R}}
\newcommand{\E}{\mathbb{E}}
\newcommand{\Prob}{\mathbb{P}}
\newcommand{\vect}[1]{\bm{#1}}
\newcommand{\mat}[1]{\bm{#1}}
\newcommand{\normtwo}[1]{\left\lVert #1 \right\rVert_2}
\newcommand{\norm}[2][]{\left\lVert #2 \right\rVert_{#1}}
\newcommand{\abs}[1]{\left\lvert #1 \right\rvert}
\newcommand{\inner}[2]{\left\langle #1,\, #2 \right\rangle}
\DeclareMathOperator*{\argmin}{arg\,min}
\DeclareMathOperator*{\argmax}{arg\,max}
\newcommand{\loss}{\mathcal{L}}
\newcommand{\dataset}{\mathcal{D}}
\newcommand{\model}{f_{\bm{\theta}}}
\newcommand{\given}{\,\vert\,}

% ---------- 図の貼り込み ----------------------------------------------
% \fig{ファイル名}{幅(\textwidth比)}  例: \fig{overview}{0.8}
\newcommand{\fig}[2]{%
  \begin{center}\includegraphics[width=#2\textwidth]{#1}\end{center}}

% \figc{ファイル名}{幅}{キャプション}
\newcommand{\figc}[3]{%
  \begin{figure}\centering
    \includegraphics[width=#2\textwidth]{#1}
    \caption{#3}
  \end{figure}}

% ---------- 表 --------------------------------------------------------
% 数値表の最良値を強調する（\best{0.812}）
\newcommand{\best}[1]{\textbf{\textcolor{mAccentBlue}{#1}}}

% ---------- 進捗報告でよく使う印 --------------------------------------
\newcommand{\done}{\textcolor{mGood}{\textbf{[完了]}}}
\newcommand{\wip}{\textcolor{mLightBrown}{\textbf{[進行中]}}}
\newcommand{\todo}{\textcolor{mMuted}{\textbf{[未着手]}}}
\newcommand{\blocked}{\textcolor{mBad}{\textbf{[停滞]}}}


%%             （common/ は .latexmkrc が TEXINPUTS に通す）
%%
%%  読み込み順には意味がある．並べ替えないこと．
%%    1. luatexja        … 和文組版の土台．最初に読む
%%    2. theme.tex       … \usetheme{metropolis}．内部で fontspec を読む
%%    3. luatexja-preset … fontspec より後に和文フォントを確定させる
%%    4. 一般パッケージ / マクロ
%% =====================================================================

% ---------- 1. 日本語（土台）------------------------------------------
\usepackage{luatexja}

% ---------- 2. テーマ（metropolis）------------------------------------
%% 2026-10-27 に ~/slides/common/ から写した（MER-DA を自己完結させるため）．以後は MER-DA/slides/common/ のこのファイルを正本とし，~/slides とは同期しない．
%% =====================================================================
%%  common/theme.tex — 見た目の設定（metropolis テーマ）
%%
%%  metropolis を選ぶ理由:
%%    - 装飾が最小限で，1枚あたりの情報量（図・数式）を最大化できる
%%    - フレームタイトルが大きく，「主張文をタイトルにする」書き方と相性が良い
%%    - プロジェクタで潰れにくい高コントラストの既定配色
%%    - 機械学習・物理系の国際会議発表で広く使われており，見慣れた形式
%% =====================================================================

\usetheme[
  progressbar=frametitle,   % フレームタイトル下に進捗バー
  numbering=fraction,       % 「7/23」形式のページ番号
  block=fill,               % block 環境に淡い背景を敷く
  sectionpage=progressbar,  % セクション扉に進捗バー
  subsectionpage=none,
  titleformat=regular,      % スモールキャップスにしない（和文と混ぜるため）
]{metropolis}

% ---------- 欧文フォント ----------------------------------------------
% metropolis 本来の指定は Fira Sans だが，本環境には未導入（警告が出るだけで
% ビルドは通る）．導入済みの Lato を代替に使う．Fira を入れた場合はこの
% ブロックを丸ごとコメントアウトすればテーマ既定に戻る．
\IfFontExistsTF{Lato}{\setsansfont{Lato}}{}
\IfFontExistsTF{DejaVu Sans Mono}{\setmonofont[Scale=0.85]{DejaVu Sans Mono}}{}

% ---------- 配色 ------------------------------------------------------
% metropolis の既定 (mDarkTeal / mLightBrown) を土台に，
% 「強調」「良い結果」「悪い結果」を区別できる色を足す．
\definecolor{mDarkTeal}{HTML}{23373B}    % 地の文・見出し
\definecolor{mLightBrown}{HTML}{EB811B}  % アクセント（alert）
\definecolor{mMediumBrown}{HTML}{C87A2F}
\definecolor{mAccentBlue}{HTML}{2A6FA8}  % 提案手法・注目対象
\definecolor{mGood}{HTML}{1B7F5C}        % 改善・良い結果
\definecolor{mBad}{HTML}{B3312B}         % 劣化・悪い結果
\definecolor{mMuted}{HTML}{7A8A8E}       % 補足・出典

% metropolis 既定の mLightBrown（オレンジ）は，白地に細い文字だと視認性が低い．
% alerted text はセクション扉と目次の文字色にも使われるので，青に置き換える．
\setbeamercolor{alerted text}{fg=mAccentBlue}
\setbeamercolor{progress bar}{fg=mAccentBlue}
\setbeamercolor{title separator}{fg=mAccentBlue}
\setbeamercolor{progress bar in head/foot}{fg=mAccentBlue}
\setbeamercolor{progress bar in section page}{fg=mAccentBlue}
% ---------- 余白・行間 ------------------------------------------------
\setbeamersize{text margin left=0.9em, text margin right=0.9em}
\setbeamertemplate{itemize items}[circle]
\setlength{\parskip}{0.4em}

% 箇条書きの階層ごとの間隔を詰める（1枚に収める調整用）
\setlength{\leftmargini}{1.2em}
\setlength{\leftmarginii}{1.0em}

% ---------- フッタに所属を出す ----------------------------------------
% \institute{} を設定していれば下部に小さく表示される
\setbeamertemplate{frame footer}{\insertshortinstitute}
\setbeamerfont{page number in head/foot}{size=\tiny}
\setbeamercolor{frame footer}{fg=mMuted}


% ---------- 3. 和文フォント -------------------------------------------
% Harano Aji は texlive-lang-japanese 同梱なので環境依存が少ない．
% Noto CJK を使う場合は [haranoaji] を [noto-otc] に差し替える．
\usepackage[haranoaji,deluxe]{luatexja-preset}
% スライドの既定は明朝ではなくゴシック（プロジェクタでの可読性）
\renewcommand{\kanjifamilydefault}{\gtdefault}
% 和欧文間の空きをやや詰める（スライドは行長が短いため）
\ltjsetparameter{xkanjiskip=0.20\zw plus 0.1\zw minus 0.05\zw}

% ---------- 4. 数式 ----------------------------------------------------
\usepackage{amsmath,amssymb}
\usepackage{mathtools}
\usepackage{bm}
\usepackage{siunitx}
\sisetup{detect-all, per-mode=symbol, group-minimum-digits=4}

% ---------- 5. 図表 ----------------------------------------------------
\usepackage{graphicx}
\usepackage{booktabs}
\usepackage{subcaption}
\usepackage{adjustbox}
\graphicspath{{figures/}{./}}

% ---------- 6. 作図 ----------------------------------------------------
\usepackage{tikz}
\usetikzlibrary{arrows.meta,positioning,calc,fit,backgrounds,shapes.geometric}
\usepackage{pgfplots}
\pgfplotsset{compat=1.18}
% pgfplots の既定を「論文の figure」ではなく「スライド」向けに寄せる
\pgfplotsset{
  every axis/.append style={
    tick label style={font=\scriptsize},
    label style={font=\small},
    legend style={font=\scriptsize, draw=none, fill=none},
    line width=0.9pt, tick style={line width=0.6pt},
  },
}

% ---------- 7. 参考文献 ------------------------------------------------
% biblatex/biber は未導入のため BibTeX + natbib を使う
\usepackage[numbers,sort&compress]{natbib}
\bibliographystyle{abbrvnat}
\setbeamertemplate{bibliography item}[text]

% ---------- 8. 体裁 ----------------------------------------------------
\usepackage{appendixnumberbeamer}   % \appendix 以降をページ番号の分母から除く
\usepackage{xcolor}
\usepackage{microtype}

% 図表キャプションは番号なし（スライド中で番号参照しないため）
\setbeamertemplate{caption}{\insertcaption}
\setbeamerfont{caption}{size=\scriptsize}

% リンク色（metropolis の配色に合わせる）
% 目次とセクション扉は内部リンクなので linkcolor が効く．
% 既定のオレンジは白地で読みにくいため，地の文と同じ色にする．
\hypersetup{colorlinks=true, allcolors=., urlcolor=mAccentBlue,
            linkcolor=mDarkTeal, citecolor=mAccentBlue}

% ---------- 9. マクロ --------------------------------------------------
%% 2026-10-27 に ~/slides/common/ から写した（MER-DA を自己完結させるため）．以後は MER-DA/slides/common/ のこのファイルを正本とし，~/slides とは同期しない．
%% =====================================================================
%%  common/macros.tex — 全スライド共通のマクロ
%%  ここに置くのは「どの発表でも使うもの」だけ．
%%  その発表固有のマクロは deck 側の main.tex に書く．
%% =====================================================================

% ---------- 強調 ------------------------------------------------------
\newcommand{\key}[1]{\textcolor{mAccentBlue}{\textbf{#1}}}  % 提案手法・注目対象
\newcommand{\good}[1]{\textcolor{mGood}{#1}}                % 改善・良い結果
\newcommand{\bad}[1]{\textcolor{mBad}{#1}}                  % 劣化・悪い結果
\newcommand{\muted}[1]{\textcolor{mMuted}{#1}}              % 補足

% 出典表示（右下に小さく．スライドの引用はこれで足りることが多い）
\newcommand{\source}[1]{%
  \vspace{-0.3em}\hfill{\tiny\color{mMuted}[#1]}\par}

% 「ひとことメッセージ」帯：そのスライドで言いたいことを1行で置く
\newcommand{\takeaway}[1]{%
  \vspace{0.4em}%
  \begin{beamercolorbox}[wd=\textwidth,sep=0.5em,rounded=true]{block body}%
    \textbf{#1}%
  \end{beamercolorbox}}

% ---------- 数式（機械学習でよく使うもの）-----------------------------
\newcommand{\R}{\mathbb{R}}
\newcommand{\E}{\mathbb{E}}
\newcommand{\Prob}{\mathbb{P}}
\newcommand{\vect}[1]{\bm{#1}}
\newcommand{\mat}[1]{\bm{#1}}
\newcommand{\normtwo}[1]{\left\lVert #1 \right\rVert_2}
\newcommand{\norm}[2][]{\left\lVert #2 \right\rVert_{#1}}
\newcommand{\abs}[1]{\left\lvert #1 \right\rvert}
\newcommand{\inner}[2]{\left\langle #1,\, #2 \right\rangle}
\DeclareMathOperator*{\argmin}{arg\,min}
\DeclareMathOperator*{\argmax}{arg\,max}
\newcommand{\loss}{\mathcal{L}}
\newcommand{\dataset}{\mathcal{D}}
\newcommand{\model}{f_{\bm{\theta}}}
\newcommand{\given}{\,\vert\,}

% ---------- 図の貼り込み ----------------------------------------------
% \fig{ファイル名}{幅(\textwidth比)}  例: \fig{overview}{0.8}
\newcommand{\fig}[2]{%
  \begin{center}\includegraphics[width=#2\textwidth]{#1}\end{center}}

% \figc{ファイル名}{幅}{キャプション}
\newcommand{\figc}[3]{%
  \begin{figure}\centering
    \includegraphics[width=#2\textwidth]{#1}
    \caption{#3}
  \end{figure}}

% ---------- 表 --------------------------------------------------------
% 数値表の最良値を強調する（\best{0.812}）
\newcommand{\best}[1]{\textbf{\textcolor{mAccentBlue}{#1}}}

% ---------- 進捗報告でよく使う印 --------------------------------------
\newcommand{\done}{\textcolor{mGood}{\textbf{[完了]}}}
\newcommand{\wip}{\textcolor{mLightBrown}{\textbf{[進行中]}}}
\newcommand{\todo}{\textcolor{mMuted}{\textbf{[未着手]}}}
\newcommand{\blocked}{\textcolor{mBad}{\textbf{[停滞]}}}

 の後に %% =====================================================================
%%  common/merda.tex — MER-DA の資料に共通の部品（上端の帯，補足のページを指す命令，色，図の部品）
%%  2026-10-27 に ~/slides/decks/MER-DA/2026-10-22-6つの問いの整理/main.tex の前置きから写した．
%%  以後はこのファイルを正本とし，~/slides とは同期しない．
%%  読み込み：各資料の main.tex で %% 2026-10-27 に ~/slides/common/ から写した（MER-DA を自己完結させるため）．以後は MER-DA/slides/common/ のこのファイルを正本とし，~/slides とは同期しない．
%% =====================================================================
%%  common/preamble.tex — 全スライド共通のプリアンブル
%%  エンジン : LuaLaTeX（日本語は luatexja）
%%  読み込み : 各 deck の main.tex から %% 2026-10-27 に ~/slides/common/ から写した（MER-DA を自己完結させるため）．以後は MER-DA/slides/common/ のこのファイルを正本とし，~/slides とは同期しない．
%% =====================================================================
%%  common/preamble.tex — 全スライド共通のプリアンブル
%%  エンジン : LuaLaTeX（日本語は luatexja）
%%  読み込み : 各 deck の main.tex から \input{preamble}
%%             （common/ は .latexmkrc が TEXINPUTS に通す）
%%
%%  読み込み順には意味がある．並べ替えないこと．
%%    1. luatexja        … 和文組版の土台．最初に読む
%%    2. theme.tex       … \usetheme{metropolis}．内部で fontspec を読む
%%    3. luatexja-preset … fontspec より後に和文フォントを確定させる
%%    4. 一般パッケージ / マクロ
%% =====================================================================

% ---------- 1. 日本語（土台）------------------------------------------
\usepackage{luatexja}

% ---------- 2. テーマ（metropolis）------------------------------------
\input{theme}

% ---------- 3. 和文フォント -------------------------------------------
% Harano Aji は texlive-lang-japanese 同梱なので環境依存が少ない．
% Noto CJK を使う場合は [haranoaji] を [noto-otc] に差し替える．
\usepackage[haranoaji,deluxe]{luatexja-preset}
% スライドの既定は明朝ではなくゴシック（プロジェクタでの可読性）
\renewcommand{\kanjifamilydefault}{\gtdefault}
% 和欧文間の空きをやや詰める（スライドは行長が短いため）
\ltjsetparameter{xkanjiskip=0.20\zw plus 0.1\zw minus 0.05\zw}

% ---------- 4. 数式 ----------------------------------------------------
\usepackage{amsmath,amssymb}
\usepackage{mathtools}
\usepackage{bm}
\usepackage{siunitx}
\sisetup{detect-all, per-mode=symbol, group-minimum-digits=4}

% ---------- 5. 図表 ----------------------------------------------------
\usepackage{graphicx}
\usepackage{booktabs}
\usepackage{subcaption}
\usepackage{adjustbox}
\graphicspath{{figures/}{./}}

% ---------- 6. 作図 ----------------------------------------------------
\usepackage{tikz}
\usetikzlibrary{arrows.meta,positioning,calc,fit,backgrounds,shapes.geometric}
\usepackage{pgfplots}
\pgfplotsset{compat=1.18}
% pgfplots の既定を「論文の figure」ではなく「スライド」向けに寄せる
\pgfplotsset{
  every axis/.append style={
    tick label style={font=\scriptsize},
    label style={font=\small},
    legend style={font=\scriptsize, draw=none, fill=none},
    line width=0.9pt, tick style={line width=0.6pt},
  },
}

% ---------- 7. 参考文献 ------------------------------------------------
% biblatex/biber は未導入のため BibTeX + natbib を使う
\usepackage[numbers,sort&compress]{natbib}
\bibliographystyle{abbrvnat}
\setbeamertemplate{bibliography item}[text]

% ---------- 8. 体裁 ----------------------------------------------------
\usepackage{appendixnumberbeamer}   % \appendix 以降をページ番号の分母から除く
\usepackage{xcolor}
\usepackage{microtype}

% 図表キャプションは番号なし（スライド中で番号参照しないため）
\setbeamertemplate{caption}{\insertcaption}
\setbeamerfont{caption}{size=\scriptsize}

% リンク色（metropolis の配色に合わせる）
% 目次とセクション扉は内部リンクなので linkcolor が効く．
% 既定のオレンジは白地で読みにくいため，地の文と同じ色にする．
\hypersetup{colorlinks=true, allcolors=., urlcolor=mAccentBlue,
            linkcolor=mDarkTeal, citecolor=mAccentBlue}

% ---------- 9. マクロ --------------------------------------------------
\input{macros}

%%             （common/ は .latexmkrc が TEXINPUTS に通す）
%%
%%  読み込み順には意味がある．並べ替えないこと．
%%    1. luatexja        … 和文組版の土台．最初に読む
%%    2. theme.tex       … \usetheme{metropolis}．内部で fontspec を読む
%%    3. luatexja-preset … fontspec より後に和文フォントを確定させる
%%    4. 一般パッケージ / マクロ
%% =====================================================================

% ---------- 1. 日本語（土台）------------------------------------------
\usepackage{luatexja}

% ---------- 2. テーマ（metropolis）------------------------------------
%% 2026-10-27 に ~/slides/common/ から写した（MER-DA を自己完結させるため）．以後は MER-DA/slides/common/ のこのファイルを正本とし，~/slides とは同期しない．
%% =====================================================================
%%  common/theme.tex — 見た目の設定（metropolis テーマ）
%%
%%  metropolis を選ぶ理由:
%%    - 装飾が最小限で，1枚あたりの情報量（図・数式）を最大化できる
%%    - フレームタイトルが大きく，「主張文をタイトルにする」書き方と相性が良い
%%    - プロジェクタで潰れにくい高コントラストの既定配色
%%    - 機械学習・物理系の国際会議発表で広く使われており，見慣れた形式
%% =====================================================================

\usetheme[
  progressbar=frametitle,   % フレームタイトル下に進捗バー
  numbering=fraction,       % 「7/23」形式のページ番号
  block=fill,               % block 環境に淡い背景を敷く
  sectionpage=progressbar,  % セクション扉に進捗バー
  subsectionpage=none,
  titleformat=regular,      % スモールキャップスにしない（和文と混ぜるため）
]{metropolis}

% ---------- 欧文フォント ----------------------------------------------
% metropolis 本来の指定は Fira Sans だが，本環境には未導入（警告が出るだけで
% ビルドは通る）．導入済みの Lato を代替に使う．Fira を入れた場合はこの
% ブロックを丸ごとコメントアウトすればテーマ既定に戻る．
\IfFontExistsTF{Lato}{\setsansfont{Lato}}{}
\IfFontExistsTF{DejaVu Sans Mono}{\setmonofont[Scale=0.85]{DejaVu Sans Mono}}{}

% ---------- 配色 ------------------------------------------------------
% metropolis の既定 (mDarkTeal / mLightBrown) を土台に，
% 「強調」「良い結果」「悪い結果」を区別できる色を足す．
\definecolor{mDarkTeal}{HTML}{23373B}    % 地の文・見出し
\definecolor{mLightBrown}{HTML}{EB811B}  % アクセント（alert）
\definecolor{mMediumBrown}{HTML}{C87A2F}
\definecolor{mAccentBlue}{HTML}{2A6FA8}  % 提案手法・注目対象
\definecolor{mGood}{HTML}{1B7F5C}        % 改善・良い結果
\definecolor{mBad}{HTML}{B3312B}         % 劣化・悪い結果
\definecolor{mMuted}{HTML}{7A8A8E}       % 補足・出典

% metropolis 既定の mLightBrown（オレンジ）は，白地に細い文字だと視認性が低い．
% alerted text はセクション扉と目次の文字色にも使われるので，青に置き換える．
\setbeamercolor{alerted text}{fg=mAccentBlue}
\setbeamercolor{progress bar}{fg=mAccentBlue}
\setbeamercolor{title separator}{fg=mAccentBlue}
\setbeamercolor{progress bar in head/foot}{fg=mAccentBlue}
\setbeamercolor{progress bar in section page}{fg=mAccentBlue}
% ---------- 余白・行間 ------------------------------------------------
\setbeamersize{text margin left=0.9em, text margin right=0.9em}
\setbeamertemplate{itemize items}[circle]
\setlength{\parskip}{0.4em}

% 箇条書きの階層ごとの間隔を詰める（1枚に収める調整用）
\setlength{\leftmargini}{1.2em}
\setlength{\leftmarginii}{1.0em}

% ---------- フッタに所属を出す ----------------------------------------
% \institute{} を設定していれば下部に小さく表示される
\setbeamertemplate{frame footer}{\insertshortinstitute}
\setbeamerfont{page number in head/foot}{size=\tiny}
\setbeamercolor{frame footer}{fg=mMuted}


% ---------- 3. 和文フォント -------------------------------------------
% Harano Aji は texlive-lang-japanese 同梱なので環境依存が少ない．
% Noto CJK を使う場合は [haranoaji] を [noto-otc] に差し替える．
\usepackage[haranoaji,deluxe]{luatexja-preset}
% スライドの既定は明朝ではなくゴシック（プロジェクタでの可読性）
\renewcommand{\kanjifamilydefault}{\gtdefault}
% 和欧文間の空きをやや詰める（スライドは行長が短いため）
\ltjsetparameter{xkanjiskip=0.20\zw plus 0.1\zw minus 0.05\zw}

% ---------- 4. 数式 ----------------------------------------------------
\usepackage{amsmath,amssymb}
\usepackage{mathtools}
\usepackage{bm}
\usepackage{siunitx}
\sisetup{detect-all, per-mode=symbol, group-minimum-digits=4}

% ---------- 5. 図表 ----------------------------------------------------
\usepackage{graphicx}
\usepackage{booktabs}
\usepackage{subcaption}
\usepackage{adjustbox}
\graphicspath{{figures/}{./}}

% ---------- 6. 作図 ----------------------------------------------------
\usepackage{tikz}
\usetikzlibrary{arrows.meta,positioning,calc,fit,backgrounds,shapes.geometric}
\usepackage{pgfplots}
\pgfplotsset{compat=1.18}
% pgfplots の既定を「論文の figure」ではなく「スライド」向けに寄せる
\pgfplotsset{
  every axis/.append style={
    tick label style={font=\scriptsize},
    label style={font=\small},
    legend style={font=\scriptsize, draw=none, fill=none},
    line width=0.9pt, tick style={line width=0.6pt},
  },
}

% ---------- 7. 参考文献 ------------------------------------------------
% biblatex/biber は未導入のため BibTeX + natbib を使う
\usepackage[numbers,sort&compress]{natbib}
\bibliographystyle{abbrvnat}
\setbeamertemplate{bibliography item}[text]

% ---------- 8. 体裁 ----------------------------------------------------
\usepackage{appendixnumberbeamer}   % \appendix 以降をページ番号の分母から除く
\usepackage{xcolor}
\usepackage{microtype}

% 図表キャプションは番号なし（スライド中で番号参照しないため）
\setbeamertemplate{caption}{\insertcaption}
\setbeamerfont{caption}{size=\scriptsize}

% リンク色（metropolis の配色に合わせる）
% 目次とセクション扉は内部リンクなので linkcolor が効く．
% 既定のオレンジは白地で読みにくいため，地の文と同じ色にする．
\hypersetup{colorlinks=true, allcolors=., urlcolor=mAccentBlue,
            linkcolor=mDarkTeal, citecolor=mAccentBlue}

% ---------- 9. マクロ --------------------------------------------------
%% 2026-10-27 に ~/slides/common/ から写した（MER-DA を自己完結させるため）．以後は MER-DA/slides/common/ のこのファイルを正本とし，~/slides とは同期しない．
%% =====================================================================
%%  common/macros.tex — 全スライド共通のマクロ
%%  ここに置くのは「どの発表でも使うもの」だけ．
%%  その発表固有のマクロは deck 側の main.tex に書く．
%% =====================================================================

% ---------- 強調 ------------------------------------------------------
\newcommand{\key}[1]{\textcolor{mAccentBlue}{\textbf{#1}}}  % 提案手法・注目対象
\newcommand{\good}[1]{\textcolor{mGood}{#1}}                % 改善・良い結果
\newcommand{\bad}[1]{\textcolor{mBad}{#1}}                  % 劣化・悪い結果
\newcommand{\muted}[1]{\textcolor{mMuted}{#1}}              % 補足

% 出典表示（右下に小さく．スライドの引用はこれで足りることが多い）
\newcommand{\source}[1]{%
  \vspace{-0.3em}\hfill{\tiny\color{mMuted}[#1]}\par}

% 「ひとことメッセージ」帯：そのスライドで言いたいことを1行で置く
\newcommand{\takeaway}[1]{%
  \vspace{0.4em}%
  \begin{beamercolorbox}[wd=\textwidth,sep=0.5em,rounded=true]{block body}%
    \textbf{#1}%
  \end{beamercolorbox}}

% ---------- 数式（機械学習でよく使うもの）-----------------------------
\newcommand{\R}{\mathbb{R}}
\newcommand{\E}{\mathbb{E}}
\newcommand{\Prob}{\mathbb{P}}
\newcommand{\vect}[1]{\bm{#1}}
\newcommand{\mat}[1]{\bm{#1}}
\newcommand{\normtwo}[1]{\left\lVert #1 \right\rVert_2}
\newcommand{\norm}[2][]{\left\lVert #2 \right\rVert_{#1}}
\newcommand{\abs}[1]{\left\lvert #1 \right\rvert}
\newcommand{\inner}[2]{\left\langle #1,\, #2 \right\rangle}
\DeclareMathOperator*{\argmin}{arg\,min}
\DeclareMathOperator*{\argmax}{arg\,max}
\newcommand{\loss}{\mathcal{L}}
\newcommand{\dataset}{\mathcal{D}}
\newcommand{\model}{f_{\bm{\theta}}}
\newcommand{\given}{\,\vert\,}

% ---------- 図の貼り込み ----------------------------------------------
% \fig{ファイル名}{幅(\textwidth比)}  例: \fig{overview}{0.8}
\newcommand{\fig}[2]{%
  \begin{center}\includegraphics[width=#2\textwidth]{#1}\end{center}}

% \figc{ファイル名}{幅}{キャプション}
\newcommand{\figc}[3]{%
  \begin{figure}\centering
    \includegraphics[width=#2\textwidth]{#1}
    \caption{#3}
  \end{figure}}

% ---------- 表 --------------------------------------------------------
% 数値表の最良値を強調する（\best{0.812}）
\newcommand{\best}[1]{\textbf{\textcolor{mAccentBlue}{#1}}}

% ---------- 進捗報告でよく使う印 --------------------------------------
\newcommand{\done}{\textcolor{mGood}{\textbf{[完了]}}}
\newcommand{\wip}{\textcolor{mLightBrown}{\textbf{[進行中]}}}
\newcommand{\todo}{\textcolor{mMuted}{\textbf{[未着手]}}}
\newcommand{\blocked}{\textcolor{mBad}{\textbf{[停滞]}}}

 の後に %% =====================================================================
%%  common/merda.tex — MER-DA の資料に共通の部品（上端の帯，補足のページを指す命令，色，図の部品）
%%  2026-10-27 に ~/slides/decks/MER-DA/2026-10-22-6つの問いの整理/main.tex の前置きから写した．
%%  以後はこのファイルを正本とし，~/slides とは同期しない．
%%  読み込み：各資料の main.tex で %% 2026-10-27 に ~/slides/common/ から写した（MER-DA を自己完結させるため）．以後は MER-DA/slides/common/ のこのファイルを正本とし，~/slides とは同期しない．
%% =====================================================================
%%  common/preamble.tex — 全スライド共通のプリアンブル
%%  エンジン : LuaLaTeX（日本語は luatexja）
%%  読み込み : 各 deck の main.tex から \input{preamble}
%%             （common/ は .latexmkrc が TEXINPUTS に通す）
%%
%%  読み込み順には意味がある．並べ替えないこと．
%%    1. luatexja        … 和文組版の土台．最初に読む
%%    2. theme.tex       … \usetheme{metropolis}．内部で fontspec を読む
%%    3. luatexja-preset … fontspec より後に和文フォントを確定させる
%%    4. 一般パッケージ / マクロ
%% =====================================================================

% ---------- 1. 日本語（土台）------------------------------------------
\usepackage{luatexja}

% ---------- 2. テーマ（metropolis）------------------------------------
\input{theme}

% ---------- 3. 和文フォント -------------------------------------------
% Harano Aji は texlive-lang-japanese 同梱なので環境依存が少ない．
% Noto CJK を使う場合は [haranoaji] を [noto-otc] に差し替える．
\usepackage[haranoaji,deluxe]{luatexja-preset}
% スライドの既定は明朝ではなくゴシック（プロジェクタでの可読性）
\renewcommand{\kanjifamilydefault}{\gtdefault}
% 和欧文間の空きをやや詰める（スライドは行長が短いため）
\ltjsetparameter{xkanjiskip=0.20\zw plus 0.1\zw minus 0.05\zw}

% ---------- 4. 数式 ----------------------------------------------------
\usepackage{amsmath,amssymb}
\usepackage{mathtools}
\usepackage{bm}
\usepackage{siunitx}
\sisetup{detect-all, per-mode=symbol, group-minimum-digits=4}

% ---------- 5. 図表 ----------------------------------------------------
\usepackage{graphicx}
\usepackage{booktabs}
\usepackage{subcaption}
\usepackage{adjustbox}
\graphicspath{{figures/}{./}}

% ---------- 6. 作図 ----------------------------------------------------
\usepackage{tikz}
\usetikzlibrary{arrows.meta,positioning,calc,fit,backgrounds,shapes.geometric}
\usepackage{pgfplots}
\pgfplotsset{compat=1.18}
% pgfplots の既定を「論文の figure」ではなく「スライド」向けに寄せる
\pgfplotsset{
  every axis/.append style={
    tick label style={font=\scriptsize},
    label style={font=\small},
    legend style={font=\scriptsize, draw=none, fill=none},
    line width=0.9pt, tick style={line width=0.6pt},
  },
}

% ---------- 7. 参考文献 ------------------------------------------------
% biblatex/biber は未導入のため BibTeX + natbib を使う
\usepackage[numbers,sort&compress]{natbib}
\bibliographystyle{abbrvnat}
\setbeamertemplate{bibliography item}[text]

% ---------- 8. 体裁 ----------------------------------------------------
\usepackage{appendixnumberbeamer}   % \appendix 以降をページ番号の分母から除く
\usepackage{xcolor}
\usepackage{microtype}

% 図表キャプションは番号なし（スライド中で番号参照しないため）
\setbeamertemplate{caption}{\insertcaption}
\setbeamerfont{caption}{size=\scriptsize}

% リンク色（metropolis の配色に合わせる）
% 目次とセクション扉は内部リンクなので linkcolor が効く．
% 既定のオレンジは白地で読みにくいため，地の文と同じ色にする．
\hypersetup{colorlinks=true, allcolors=., urlcolor=mAccentBlue,
            linkcolor=mDarkTeal, citecolor=mAccentBlue}

% ---------- 9. マクロ --------------------------------------------------
\input{macros}
 の後に %% =====================================================================
%%  common/merda.tex — MER-DA の資料に共通の部品（上端の帯，補足のページを指す命令，色，図の部品）
%%  2026-10-27 に ~/slides/decks/MER-DA/2026-10-22-6つの問いの整理/main.tex の前置きから写した．
%%  以後はこのファイルを正本とし，~/slides とは同期しない．
%%  読み込み：各資料の main.tex で \input{preamble} の後に \input{merda}
%% =====================================================================
\usepackage{multirow}
\usepackage{array}
\usepackage{colortbl}
\usepackage{tabularx}
\usepgfplotslibrary{groupplots}
\pgfplotsset{legend cell align=left}
\usetikzlibrary{decorations.pathreplacing,patterns}
\metroset{sectionpage=none}

\setbeamertemplate{frame numbering}{\insertpagenumber/\insertdocumentendpage}
\setbeamerfont{page number in head/foot}{size=\tiny}

% ---------- 上端の帯：部の名前｜スライドの役割 ----------------------------
%   \setband{部の名前}   部の最初のスライドの前で 1 回
%   \rl{役割}            各スライドの本文の最初の行（目的・方法・結果・まとめなど）
\newcommand{\bandq}{}\newcommand{\bandrole}{}
\newcommand{\setband}[1]{\gdef\bandq{#1}}
\newcommand{\rl}[1]{\gdef\bandrole{#1}}
\definecolor{bandbg}{HTML}{E9EEF0}
\newcommand{\bandheadline}{\setbeamertemplate{headline}{%
  \ifx\bandq\empty\else
  \begin{beamercolorbox}[wd=\paperwidth,ht=1.9ex,dp=0.7ex,leftskip=0.9em,rightskip=0.9em]{headline}%
    \colorbox{bandbg}{\makebox[\dimexpr\paperwidth-1.8em-2\fboxsep\relax][l]{\tiny\color{mDarkTeal}%
      \textbf{\bandq}\ifx\bandrole\empty\else\enspace{}｜\enspace{}\textbf{\bandrole}\fi}}%
  \end{beamercolorbox}\fi}}
\bandheadline   % metropolis は \appendix で headline を消すので，\appendix の後でもう一度呼ぶ

% ---------- 本編のスライドから補足のページを指す ------------------------------
%   \hosoku{ラベル}{何があるか}：左下のフッターに「補足 p.N：…」と出し，次のスライドには出さない
\def\hosokutext{}
\newcommand{\hosoku}[2]{\gdef\hosokutext{\textcolor{mDarkTeal}{補足 p.~\pageref{#1}：#2}}}
\setbeamertemplate{frame footer}{\ifx\hosokutext\empty\insertshortinstitute\else\hosokutext\fi\global\let\hosokutext\empty}

% ---------- 箱と帯 ----------------------------------------------------------
\renewcommand{\source}[1]{\par\vspace{-0.45em}\hfill{\tiny\color{mMuted}[#1]}\par}
\renewcommand{\takeaway}[1]{\vspace{0.12em}\begin{beamercolorbox}[wd=\textwidth,sep=0.34em,rounded=true]{block body}\textbf{#1}\end{beamercolorbox}}
\newcommand{\band}[3]{\par\noindent\fcolorbox{#1}{#1!7}{\parbox{\dimexpr\linewidth-2\fboxsep-2\fboxrule}{\raggedright\small\textbf{\textcolor{#1}{#2}}\enspace #3}}\par}
\definecolor{cAsk}{HTML}{5B4A9E}\definecolor{cHyp}{HTML}{8A6418}
\newcommand{\toi}[1]{\band{cAsk}{問い}{#1}}
\newcommand{\imi}[1]{\band{mDarkTeal}{この結果の意味}{#1}}
\newcommand{\tighteq}{\setlength{\abovedisplayskip}{3pt}\setlength{\belowdisplayskip}{3pt}\setlength{\abovedisplayshortskip}{2pt}\setlength{\belowdisplayshortskip}{2pt}}
\newcommand{\crit}[1]{\textcolor{cCrit}{\textbf{#1}}}
\newcommand{\disc}[1]{\textcolor{cDisc}{\textbf{#1}}}
\ltjsetparameter{prebreakpenalty={`〜,10000}}\ltjsetparameter{postbreakpenalty={`〜,10000}}\ltjsetparameter{prebreakpenalty={`／,10000}}   % 範囲の「〜」の前後と，並べた値の「／」の前で改行しない

% ---------- 色 ----------------------------------------------------------------
\definecolor{cCrit}{HTML}{D9822B}   % 判断基準のずれ o・切片
\definecolor{cDisc}{HTML}{2A6FA8}   % 識別力 d・重みの向き
\definecolor{cSrc}{HTML}{D9E6F2}    % ソース
\definecolor{cTgt}{HTML}{F6E3C8}    % ターゲット
\definecolor{cMod}{HTML}{E4D9EF}    % 別のモダリティ

% ---------- 図の部品 ----------------------------------------------------------
\pgfplotsset{
  table/search path={figures},
  table/col sep=comma,
  table/comment chars={\#},
  every axis/.append style={label style={font=\scriptsize}, tick label style={font=\scriptsize}},
}
\tikzset{
  srcbox/.style={draw=mDarkTeal!70, fill=cSrc, rounded corners=2pt, align=center, font=\scriptsize, minimum height=0.6cm},
  tgtbox/.style={draw=mDarkTeal!70, fill=cTgt, rounded corners=2pt, align=center, font=\scriptsize, minimum height=0.6cm},
  modbox/.style={draw=mDarkTeal!70, fill=cMod, rounded corners=2pt, align=center, font=\scriptsize, minimum height=0.6cm},
  clf/.style={draw=mDarkTeal, fill=white, rounded corners=2pt, align=center, font=\scriptsize, minimum height=0.55cm},
  arr/.style={-{Stealth[length=2.2mm]}, line width=0.8pt, mDarkTeal},
  testbox/.style={draw=mBad, densely dashed, line width=0.8pt, fill=mBad!7, rounded corners=2pt, align=center, font=\scriptsize, minimum height=0.6cm},
  encbox/.style={draw=mMuted, fill=mMuted!22, rounded corners=2pt, align=center, font=\scriptsize},
  tarr/.style={-{Stealth[length=2.2mm]}, line width=0.8pt, mBad, densely dashed},
  warr/.style={-{Stealth[length=2.2mm]}, line width=0.8pt, cHyp, densely dotted},
  tag/.style={rounded corners=1.5pt, inner sep=1.5pt, font=\tiny\bfseries, text=white},
  qbox/.style={draw=mDarkTeal!70, fill=white, rounded corners=2pt, align=left, font=\scriptsize, inner sep=3pt},
}

%% =====================================================================
\usepackage{multirow}
\usepackage{array}
\usepackage{colortbl}
\usepackage{tabularx}
\usepgfplotslibrary{groupplots}
\pgfplotsset{legend cell align=left}
\usetikzlibrary{decorations.pathreplacing,patterns}
\metroset{sectionpage=none}

\setbeamertemplate{frame numbering}{\insertpagenumber/\insertdocumentendpage}
\setbeamerfont{page number in head/foot}{size=\tiny}

% ---------- 上端の帯：部の名前｜スライドの役割 ----------------------------
%   \setband{部の名前}   部の最初のスライドの前で 1 回
%   \rl{役割}            各スライドの本文の最初の行（目的・方法・結果・まとめなど）
\newcommand{\bandq}{}\newcommand{\bandrole}{}
\newcommand{\setband}[1]{\gdef\bandq{#1}}
\newcommand{\rl}[1]{\gdef\bandrole{#1}}
\definecolor{bandbg}{HTML}{E9EEF0}
\newcommand{\bandheadline}{\setbeamertemplate{headline}{%
  \ifx\bandq\empty\else
  \begin{beamercolorbox}[wd=\paperwidth,ht=1.9ex,dp=0.7ex,leftskip=0.9em,rightskip=0.9em]{headline}%
    \colorbox{bandbg}{\makebox[\dimexpr\paperwidth-1.8em-2\fboxsep\relax][l]{\tiny\color{mDarkTeal}%
      \textbf{\bandq}\ifx\bandrole\empty\else\enspace{}｜\enspace{}\textbf{\bandrole}\fi}}%
  \end{beamercolorbox}\fi}}
\bandheadline   % metropolis は \appendix で headline を消すので，\appendix の後でもう一度呼ぶ

% ---------- 本編のスライドから補足のページを指す ------------------------------
%   \hosoku{ラベル}{何があるか}：左下のフッターに「補足 p.N：…」と出し，次のスライドには出さない
\def\hosokutext{}
\newcommand{\hosoku}[2]{\gdef\hosokutext{\textcolor{mDarkTeal}{補足 p.~\pageref{#1}：#2}}}
\setbeamertemplate{frame footer}{\ifx\hosokutext\empty\insertshortinstitute\else\hosokutext\fi\global\let\hosokutext\empty}

% ---------- 箱と帯 ----------------------------------------------------------
\renewcommand{\source}[1]{\par\vspace{-0.45em}\hfill{\tiny\color{mMuted}[#1]}\par}
\renewcommand{\takeaway}[1]{\vspace{0.12em}\begin{beamercolorbox}[wd=\textwidth,sep=0.34em,rounded=true]{block body}\textbf{#1}\end{beamercolorbox}}
\newcommand{\band}[3]{\par\noindent\fcolorbox{#1}{#1!7}{\parbox{\dimexpr\linewidth-2\fboxsep-2\fboxrule}{\raggedright\small\textbf{\textcolor{#1}{#2}}\enspace #3}}\par}
\definecolor{cAsk}{HTML}{5B4A9E}\definecolor{cHyp}{HTML}{8A6418}
\newcommand{\toi}[1]{\band{cAsk}{問い}{#1}}
\newcommand{\imi}[1]{\band{mDarkTeal}{この結果の意味}{#1}}
\newcommand{\tighteq}{\setlength{\abovedisplayskip}{3pt}\setlength{\belowdisplayskip}{3pt}\setlength{\abovedisplayshortskip}{2pt}\setlength{\belowdisplayshortskip}{2pt}}
\newcommand{\crit}[1]{\textcolor{cCrit}{\textbf{#1}}}
\newcommand{\disc}[1]{\textcolor{cDisc}{\textbf{#1}}}
\ltjsetparameter{prebreakpenalty={`〜,10000}}\ltjsetparameter{postbreakpenalty={`〜,10000}}\ltjsetparameter{prebreakpenalty={`／,10000}}   % 範囲の「〜」の前後と，並べた値の「／」の前で改行しない

% ---------- 色 ----------------------------------------------------------------
\definecolor{cCrit}{HTML}{D9822B}   % 判断基準のずれ o・切片
\definecolor{cDisc}{HTML}{2A6FA8}   % 識別力 d・重みの向き
\definecolor{cSrc}{HTML}{D9E6F2}    % ソース
\definecolor{cTgt}{HTML}{F6E3C8}    % ターゲット
\definecolor{cMod}{HTML}{E4D9EF}    % 別のモダリティ

% ---------- 図の部品 ----------------------------------------------------------
\pgfplotsset{
  table/search path={figures},
  table/col sep=comma,
  table/comment chars={\#},
  every axis/.append style={label style={font=\scriptsize}, tick label style={font=\scriptsize}},
}
\tikzset{
  srcbox/.style={draw=mDarkTeal!70, fill=cSrc, rounded corners=2pt, align=center, font=\scriptsize, minimum height=0.6cm},
  tgtbox/.style={draw=mDarkTeal!70, fill=cTgt, rounded corners=2pt, align=center, font=\scriptsize, minimum height=0.6cm},
  modbox/.style={draw=mDarkTeal!70, fill=cMod, rounded corners=2pt, align=center, font=\scriptsize, minimum height=0.6cm},
  clf/.style={draw=mDarkTeal, fill=white, rounded corners=2pt, align=center, font=\scriptsize, minimum height=0.55cm},
  arr/.style={-{Stealth[length=2.2mm]}, line width=0.8pt, mDarkTeal},
  testbox/.style={draw=mBad, densely dashed, line width=0.8pt, fill=mBad!7, rounded corners=2pt, align=center, font=\scriptsize, minimum height=0.6cm},
  encbox/.style={draw=mMuted, fill=mMuted!22, rounded corners=2pt, align=center, font=\scriptsize},
  tarr/.style={-{Stealth[length=2.2mm]}, line width=0.8pt, mBad, densely dashed},
  warr/.style={-{Stealth[length=2.2mm]}, line width=0.8pt, cHyp, densely dotted},
  tag/.style={rounded corners=1.5pt, inner sep=1.5pt, font=\tiny\bfseries, text=white},
  qbox/.style={draw=mDarkTeal!70, fill=white, rounded corners=2pt, align=left, font=\scriptsize, inner sep=3pt},
}

%% =====================================================================
\usepackage{multirow}
\usepackage{array}
\usepackage{colortbl}
\usepackage{tabularx}
\usepgfplotslibrary{groupplots}
\pgfplotsset{legend cell align=left}
\usetikzlibrary{decorations.pathreplacing,patterns}
\metroset{sectionpage=none}

\setbeamertemplate{frame numbering}{\insertpagenumber/\insertdocumentendpage}
\setbeamerfont{page number in head/foot}{size=\tiny}

% ---------- 上端の帯：部の名前｜スライドの役割 ----------------------------
%   \setband{部の名前}   部の最初のスライドの前で 1 回
%   \rl{役割}            各スライドの本文の最初の行（目的・方法・結果・まとめなど）
\newcommand{\bandq}{}\newcommand{\bandrole}{}
\newcommand{\setband}[1]{\gdef\bandq{#1}}
\newcommand{\rl}[1]{\gdef\bandrole{#1}}
\definecolor{bandbg}{HTML}{E9EEF0}
\newcommand{\bandheadline}{\setbeamertemplate{headline}{%
  \ifx\bandq\empty\else
  \begin{beamercolorbox}[wd=\paperwidth,ht=1.9ex,dp=0.7ex,leftskip=0.9em,rightskip=0.9em]{headline}%
    \colorbox{bandbg}{\makebox[\dimexpr\paperwidth-1.8em-2\fboxsep\relax][l]{\tiny\color{mDarkTeal}%
      \textbf{\bandq}\ifx\bandrole\empty\else\enspace{}｜\enspace{}\textbf{\bandrole}\fi}}%
  \end{beamercolorbox}\fi}}
\bandheadline   % metropolis は \appendix で headline を消すので，\appendix の後でもう一度呼ぶ

% ---------- 本編のスライドから補足のページを指す ------------------------------
%   \hosoku{ラベル}{何があるか}：左下のフッターに「補足 p.N：…」と出し，次のスライドには出さない
\def\hosokutext{}
\newcommand{\hosoku}[2]{\gdef\hosokutext{\textcolor{mDarkTeal}{補足 p.~\pageref{#1}：#2}}}
\setbeamertemplate{frame footer}{\ifx\hosokutext\empty\insertshortinstitute\else\hosokutext\fi\global\let\hosokutext\empty}

% ---------- 箱と帯 ----------------------------------------------------------
\renewcommand{\source}[1]{\par\vspace{-0.45em}\hfill{\tiny\color{mMuted}[#1]}\par}
\renewcommand{\takeaway}[1]{\vspace{0.12em}\begin{beamercolorbox}[wd=\textwidth,sep=0.34em,rounded=true]{block body}\textbf{#1}\end{beamercolorbox}}
\newcommand{\band}[3]{\par\noindent\fcolorbox{#1}{#1!7}{\parbox{\dimexpr\linewidth-2\fboxsep-2\fboxrule}{\raggedright\small\textbf{\textcolor{#1}{#2}}\enspace #3}}\par}
\definecolor{cAsk}{HTML}{5B4A9E}\definecolor{cHyp}{HTML}{8A6418}
\newcommand{\toi}[1]{\band{cAsk}{問い}{#1}}
\newcommand{\imi}[1]{\band{mDarkTeal}{この結果の意味}{#1}}
\newcommand{\tighteq}{\setlength{\abovedisplayskip}{3pt}\setlength{\belowdisplayskip}{3pt}\setlength{\abovedisplayshortskip}{2pt}\setlength{\belowdisplayshortskip}{2pt}}
\newcommand{\crit}[1]{\textcolor{cCrit}{\textbf{#1}}}
\newcommand{\disc}[1]{\textcolor{cDisc}{\textbf{#1}}}
\ltjsetparameter{prebreakpenalty={`〜,10000}}\ltjsetparameter{postbreakpenalty={`〜,10000}}\ltjsetparameter{prebreakpenalty={`／,10000}}   % 範囲の「〜」の前後と，並べた値の「／」の前で改行しない

% ---------- 色 ----------------------------------------------------------------
\definecolor{cCrit}{HTML}{D9822B}   % 判断基準のずれ o・切片
\definecolor{cDisc}{HTML}{2A6FA8}   % 識別力 d・重みの向き
\definecolor{cSrc}{HTML}{D9E6F2}    % ソース
\definecolor{cTgt}{HTML}{F6E3C8}    % ターゲット
\definecolor{cMod}{HTML}{E4D9EF}    % 別のモダリティ

% ---------- 図の部品 ----------------------------------------------------------
\pgfplotsset{
  table/search path={figures},
  table/col sep=comma,
  table/comment chars={\#},
  every axis/.append style={label style={font=\scriptsize}, tick label style={font=\scriptsize}},
}
\tikzset{
  srcbox/.style={draw=mDarkTeal!70, fill=cSrc, rounded corners=2pt, align=center, font=\scriptsize, minimum height=0.6cm},
  tgtbox/.style={draw=mDarkTeal!70, fill=cTgt, rounded corners=2pt, align=center, font=\scriptsize, minimum height=0.6cm},
  modbox/.style={draw=mDarkTeal!70, fill=cMod, rounded corners=2pt, align=center, font=\scriptsize, minimum height=0.6cm},
  clf/.style={draw=mDarkTeal, fill=white, rounded corners=2pt, align=center, font=\scriptsize, minimum height=0.55cm},
  arr/.style={-{Stealth[length=2.2mm]}, line width=0.8pt, mDarkTeal},
  testbox/.style={draw=mBad, densely dashed, line width=0.8pt, fill=mBad!7, rounded corners=2pt, align=center, font=\scriptsize, minimum height=0.6cm},
  encbox/.style={draw=mMuted, fill=mMuted!22, rounded corners=2pt, align=center, font=\scriptsize},
  tarr/.style={-{Stealth[length=2.2mm]}, line width=0.8pt, mBad, densely dashed},
  warr/.style={-{Stealth[length=2.2mm]}, line width=0.8pt, cHyp, densely dotted},
  tag/.style={rounded corners=1.5pt, inner sep=1.5pt, font=\tiny\bfseries, text=white},
  qbox/.style={draw=mDarkTeal!70, fill=white, rounded corners=2pt, align=left, font=\scriptsize, inner sep=3pt},
}

%% =====================================================================
\usepackage{multirow}
\usepackage{array}
\usepackage{colortbl}
\usepackage{tabularx}
\usepgfplotslibrary{groupplots}
\pgfplotsset{legend cell align=left}
\usetikzlibrary{decorations.pathreplacing,patterns}
\metroset{sectionpage=none}

\setbeamertemplate{frame numbering}{\insertpagenumber/\insertdocumentendpage}
\setbeamerfont{page number in head/foot}{size=\tiny}

% ---------- 上端の帯：部の名前｜スライドの役割 ----------------------------
%   \setband{部の名前}   部の最初のスライドの前で 1 回
%   \rl{役割}            各スライドの本文の最初の行（目的・方法・結果・まとめなど）
\newcommand{\bandq}{}\newcommand{\bandrole}{}
\newcommand{\setband}[1]{\gdef\bandq{#1}}
\newcommand{\rl}[1]{\gdef\bandrole{#1}}
\definecolor{bandbg}{HTML}{E9EEF0}
\newcommand{\bandheadline}{\setbeamertemplate{headline}{%
  \ifx\bandq\empty\else
  \begin{beamercolorbox}[wd=\paperwidth,ht=1.9ex,dp=0.7ex,leftskip=0.9em,rightskip=0.9em]{headline}%
    \colorbox{bandbg}{\makebox[\dimexpr\paperwidth-1.8em-2\fboxsep\relax][l]{\tiny\color{mDarkTeal}%
      \textbf{\bandq}\ifx\bandrole\empty\else\enspace{}｜\enspace{}\textbf{\bandrole}\fi}}%
  \end{beamercolorbox}\fi}}
\bandheadline   % metropolis は \appendix で headline を消すので，\appendix の後でもう一度呼ぶ

% ---------- 本編のスライドから補足のページを指す ------------------------------
%   \hosoku{ラベル}{何があるか}：左下のフッターに「補足 p.N：…」と出し，次のスライドには出さない
\def\hosokutext{}
\newcommand{\hosoku}[2]{\gdef\hosokutext{\textcolor{mDarkTeal}{補足 p.~\pageref{#1}：#2}}}
\setbeamertemplate{frame footer}{\ifx\hosokutext\empty\insertshortinstitute\else\hosokutext\fi\global\let\hosokutext\empty}

% ---------- 箱と帯 ----------------------------------------------------------
\renewcommand{\source}[1]{\par\vspace{-0.45em}\hfill{\tiny\color{mMuted}[#1]}\par}
\renewcommand{\takeaway}[1]{\vspace{0.12em}\begin{beamercolorbox}[wd=\textwidth,sep=0.34em,rounded=true]{block body}\textbf{#1}\end{beamercolorbox}}
\newcommand{\band}[3]{\par\noindent\fcolorbox{#1}{#1!7}{\parbox{\dimexpr\linewidth-2\fboxsep-2\fboxrule}{\raggedright\small\textbf{\textcolor{#1}{#2}}\enspace #3}}\par}
\definecolor{cAsk}{HTML}{5B4A9E}\definecolor{cHyp}{HTML}{8A6418}
\newcommand{\toi}[1]{\band{cAsk}{問い}{#1}}
\newcommand{\imi}[1]{\band{mDarkTeal}{この結果の意味}{#1}}
\newcommand{\tighteq}{\setlength{\abovedisplayskip}{3pt}\setlength{\belowdisplayskip}{3pt}\setlength{\abovedisplayshortskip}{2pt}\setlength{\belowdisplayshortskip}{2pt}}
\newcommand{\crit}[1]{\textcolor{cCrit}{\textbf{#1}}}
\newcommand{\disc}[1]{\textcolor{cDisc}{\textbf{#1}}}
\ltjsetparameter{prebreakpenalty={`〜,10000}}\ltjsetparameter{postbreakpenalty={`〜,10000}}\ltjsetparameter{prebreakpenalty={`／,10000}}   % 範囲の「〜」の前後と，並べた値の「／」の前で改行しない

% ---------- 色 ----------------------------------------------------------------
\definecolor{cCrit}{HTML}{D9822B}   % 判断基準のずれ o・切片
\definecolor{cDisc}{HTML}{2A6FA8}   % 識別力 d・重みの向き
\definecolor{cSrc}{HTML}{D9E6F2}    % ソース
\definecolor{cTgt}{HTML}{F6E3C8}    % ターゲット
\definecolor{cMod}{HTML}{E4D9EF}    % 別のモダリティ

% ---------- 図の部品 ----------------------------------------------------------
\pgfplotsset{
  table/search path={figures},
  table/col sep=comma,
  table/comment chars={\#},
  every axis/.append style={label style={font=\scriptsize}, tick label style={font=\scriptsize}},
}
\tikzset{
  srcbox/.style={draw=mDarkTeal!70, fill=cSrc, rounded corners=2pt, align=center, font=\scriptsize, minimum height=0.6cm},
  tgtbox/.style={draw=mDarkTeal!70, fill=cTgt, rounded corners=2pt, align=center, font=\scriptsize, minimum height=0.6cm},
  modbox/.style={draw=mDarkTeal!70, fill=cMod, rounded corners=2pt, align=center, font=\scriptsize, minimum height=0.6cm},
  clf/.style={draw=mDarkTeal, fill=white, rounded corners=2pt, align=center, font=\scriptsize, minimum height=0.55cm},
  arr/.style={-{Stealth[length=2.2mm]}, line width=0.8pt, mDarkTeal},
  testbox/.style={draw=mBad, densely dashed, line width=0.8pt, fill=mBad!7, rounded corners=2pt, align=center, font=\scriptsize, minimum height=0.6cm},
  encbox/.style={draw=mMuted, fill=mMuted!22, rounded corners=2pt, align=center, font=\scriptsize},
  tarr/.style={-{Stealth[length=2.2mm]}, line width=0.8pt, mBad, densely dashed},
  warr/.style={-{Stealth[length=2.2mm]}, line width=0.8pt, cHyp, densely dotted},
  tag/.style={rounded corners=1.5pt, inner sep=1.5pt, font=\tiny\bfseries, text=white},
  qbox/.style={draw=mDarkTeal!70, fill=white, rounded corners=2pt, align=left, font=\scriptsize, inner sep=3pt},
}
